\section{Introduction}

Editing a real image according to a text instruction is deceptively simple to describe yet remarkably difficult to achieve: the model must simultaneously understand what to change, what to leave intact, and how to reconcile the two, all without retraining on the target image. Training-free approaches are particularly attractive for this task, as they repurpose pre-trained generative models at inference time, avoiding the prohibitive cost of per-image optimization.
Existing training-free editors predominantly build on diffusion models, where editing is realized by manipulating the denoising trajectory of an inverted image. Prompt-to-Prompt~\cite{prompt2prompt} and its successors~\cite{Null-text,PnP,PnP_Inversion,masactrl,pix2pix-zero,LEDits++} demonstrate that cross-attention and self-attention manipulations during denoising can localize edits to semantically relevant regions. More recently, rectified flow models~\cite{Flowedit,RF-inversion,ReFlex,FireFlow} offer near-linear ODE trajectories that simplify inversion and reduce reconstruction error. Despite their diversity, all these methods share a common structural dependency: editing a real image requires first inverting it into the model's noise space, then re-generating under a modified prompt. This creates two persistent failure modes. First, \textit{inversion error}: even small discrepancies in the reconstructed trajectory propagate through every subsequent generation step, producing visual artifacts and semantic drift. Second, \textit{global feature entanglement}: because text conditioning modulates the entire denoising trajectory at every timestep, a local edit unavoidably perturbs the global activation pattern, causing unintended changes in structurally unrelated regions.
These challenges are not incidental implementation details; they are structural consequences of operating in a continuous noise space where local and global information are inextricably coupled. This motivates a natural question: \textit{is there a generative paradigm where such coupling is absent by design?}
Visual Autoregressive Modeling (VAR)~\cite{VAR} offers an affirmative answer. Rather than denoising a continuous noise map, VAR generates images by predicting discrete token sequences at progressively finer scales, akin to progressively filling in a sketch with increasing detail. Crucially, each token corresponds to a well-defined spatial patch at a specific scale, and generation proceeds by sampling one token at a time. Editing is therefore, in principle, as simple as \textit{selectively rewriting the tokens that correspond to the target region}, with no inversion, no trajectory, and no global coupling to untangle. The remaining challenge is knowing \textit{which tokens to rewrite}: a question of spatial understanding, not generative mechanics. AREdit~\cite{AREdit} took a first step in this direction by caching token statistics from the source generation pass, but its design treats token selection as a calibration problem rather than a spatial understanding problem; consequently, it cannot generalize across editing types without per-category thresholds supplied by the user, transferring the burden of spatial reasoning back to manual tuning rather than resolving it.
We answer the spatial understanding question by exploiting the model's own cross-modal attention maps between text and visual tokens. At each generation scale, we aggregate attention weights across transformer layers and apply IQR-based filtering to identify and suppress layers whose attention distributions are dominated by positional bias rather than semantic content, a statistically principled criterion that requires no manual layer selection. The filtered maps are then binarized to obtain a spatial editing mask: tokens falling within the mask are resampled under the target prompt, while all others are directly transplanted from the source image's discrete BSQ-VAE~\cite{BSQVAE} token encoding, a lossless operation that guarantees pixel-faithful preservation of unedited regions by construction. When a reference mask is available, a binary search procedure dynamically calibrates the binarization threshold to the spatial distribution of each instance. This is motivated by a straightforward observation: editing-region extent varies substantially across instances, and any fixed threshold implicitly assumes a uniform spatial distribution, an assumption that open-vocabulary instructions routinely violate. Together, these components realize an \textit{understanding-before-editing} pipeline: the model first resolves spatial semantics from its own attention, then rewrites only what the language instruction specifies, with no inversion, no continuous feature injection, and no per-category tuning.

Our contributions are as follows:

\begin{itemize}
    \item We propose the first attention-map-driven editing framework for VAR-based autoregressive models, achieving localized editing through selective discrete token replacement at each generation scale, a mechanism that is structurally inversion-free and free of global trajectory coupling by construction.

    \item We introduce IQR-filtered attention aggregation and binary-search mask thresholding, an adaptive masking strategy that suppresses layer-level positional artifacts and automatically calibrates spatial precision to each editing instance, requiring no external segmentation, user-provided masks, or per-category hyperparameter tuning.

    \item We extend the framework to real image editing by directly encoding the source image as BSQ-VAE tokens and transplanting them into unedited regions, a lossless operation that guarantees pixel-faithful background preservation by construction without continuous feature injection or reconstruction distortion. Evaluated on PIE-Bench, a 2B-parameter VAR model under our framework outperforms state-of-the-art diffusion-based (2B) and flow-based (12B) baselines across nearly all metrics, completing editing in 2.11 seconds, the fastest among all compared approaches.
\end{itemize}
\input{imgs/pipeline}